% Paper analysis section of the A2.2 report (SegFormer, Xie et al., NeurIPS 2021).
% Every number about the paper below was checked against the arXiv version (2105.15203).
\section{Phân tích paper SegFormer}\label{sec:paper}

\subsection{Thông tin và lý do chọn}
E.~Xie, W.~Wang, Z.~Yu, A.~Anandkumar, J.~M.~Alvarez, P.~Luo. \emph{SegFormer: Simple and Efficient Design for Semantic
Segmentation with Transformers}. NeurIPS 2021~\cite{xie2021}. Mã nguồn chính thức: \url{https://github.com/NVlabs/SegFormer};
trọng số encoder dùng trong bài: \code{nvidia/mit-b2} trên Hugging Face~\cite{wolf2020}. Nhóm chọn paper này vì một trong hai
mô hình của nhóm chính là SegFormer-B2, nên các nhận định của paper kiểm chứng được trực tiếp bằng thí nghiệm của nhóm.

\subsection{Vấn đề paper giải quyết}
Trước SegFormer, các mô hình phân đoạn dùng Transformer gặp hai trở ngại. Thứ nhất, ViT gốc chỉ cho đặc trưng ở một tỉ lệ
và độ phân giải thấp, trong khi phân đoạn cần cả chi tiết nhỏ (biên vật) lẫn ngữ cảnh rộng. Thứ hai, ViT dùng
positional encoding có kích thước cố định; khi ảnh test lớn hơn ảnh train, positional encoding phải nội suy và độ chính
xác giảm. Các hướng đi trước đó (SETR, PVT, Swin) chủ yếu thiết kế lại encoder và vẫn dùng bộ giải mã nặng. SegFormer
đặt giả thiết: \st{nếu encoder đủ tốt ở nhiều tỉ lệ thì bộ giải mã có thể rất nhẹ}.

\subsection{Phương pháp}
\textbf{Encoder Mix Transformer (MiT).} Gồm 4 tầng cho ra đặc trưng ở 1/4, 1/8, 1/16, 1/32 độ phân giải ảnh, giống kim
tự tháp đặc trưng của CNN. Mỗi tầng có ba điểm khác ViT:
\begin{itemize}
  \item \st{Overlapped patch merging:} cắt patch bằng tích chập có bước trượt nhỏ hơn kích thước nhân (ví dụ nhân 7, bước 4
        ở tầng đầu), nên các patch chồng lên nhau và giữ được tính liên tục cục bộ quanh biên patch.
  \item \st{Efficient self-attention:} trước khi tính attention, chuỗi key/value được rút ngắn $R$ lần ($R = 64, 16, 4, 1$
        ở 4 tầng, tức giảm 8, 4, 2, 1 lần theo mỗi chiều không gian), giảm chi phí từ $O(N^2)$ xuống $O(N^2/R)$. Nhờ vậy attention toàn cục vẫn
        chạy được ở độ phân giải 1/4.
  \item \st{Mix-FFN:} chèn một tích chập depthwise $3\times3$ vào khối FFN. Tích chập có zero-padding ``rò rỉ'' thông tin
        vị trí, nên mô hình \st{không cần positional encoding} và không bị ràng buộc vào kích thước ảnh lúc train.
\end{itemize}
\textbf{Bộ giải mã toàn MLP.} Bốn bản đồ đặc trưng được chiếu về cùng $C$ kênh bằng một lớp tuyến tính, phóng lên 1/4 độ
phân giải, ghép lại, trộn bằng một lớp tuyến tính nữa rồi dự đoán $N_{cls}$ lớp. Với B2, $C = 768$ và bộ giải mã chỉ có
3,3 triệu tham số. Paper lập luận rằng tầng nông của MiT đã nhìn cục bộ còn tầng sâu nhìn gần như toàn ảnh (minh hoạ
bằng effective receptive field), nên ghép đặc trưng các tầng là đủ, không cần module ngữ cảnh như ASPP.

\textbf{Loss và huấn luyện.} Chỉ dùng cross-entropy theo điểm ảnh; paper nói rõ không dùng OHEM, loss phụ hay cân bằng lớp.
AdamW, learning rate $6\times10^{-5}$, lịch ``poly'' hệ số 1,0; encoder pretrained ImageNet-1k, bộ giải mã khởi tạo ngẫu
nhiên. Augmentation: thay đổi kích thước ngẫu nhiên tỉ lệ 0,5--2,0, lật ngang, cắt ngẫu nhiên $512\times512$ (ADE20K,
COCO-Stuff) hoặc $1024\times1024$ (Cityscapes); 160 nghìn vòng lặp, batch 16, trên 8 GPU V100.

\subsection{Thí nghiệm của paper}
Ba tập: ADE20K (150 lớp, 20.210 ảnh), Cityscapes (19 lớp, 5.000 ảnh độ phân giải cao) và COCO-Stuff (172 lớp, 164 nghìn
ảnh); metric mIoU, kèm số tham số, FLOPs và tốc độ. Các kết quả nhóm thấy quan trọng nhất (Bảng~\ref{tab:paper-numbers}):

\begin{table}[H]
\centering
\caption{Một số kết quả của paper được nhóm đối chiếu (số lấy từ Bảng 1 và Bảng 5 của paper).}\label{tab:paper-numbers}
\begin{tabularx}{\linewidth}{@{}P{4.6cm}L@{}}
\toprule
Thí nghiệm & Kết quả \\\midrule
Cỡ mô hình (Bảng 1a) & SegFormer-B2: encoder 24,2M + decoder 3,3M tham số; mIoU 46,5 (ADE20K), 81,0 (Cityscapes), 44,6 (COCO-Stuff), test một tỉ lệ \\
Bộ giải mã MLP với encoder CNN (Bảng 1d, ADE20K) & ResNet-50 + bộ giải mã MLP: 34,7 mIoU; MiT-B2 + cùng bộ giải mã: 45,4. Bộ giải mã nhẹ chỉ hiệu quả khi encoder là Transformer \\
Mix-FFN và positional encoding (Bảng 1c, Cityscapes) & Train $768\times768$, test $1024\times2048$: dùng positional encoding giảm 77,3 $\rightarrow$ 74,0; dùng Mix-FFN chỉ giảm 80,5 $\rightarrow$ 79,8 \\
Số kênh $C$ của bộ giải mã (Bảng 1b) & 256 $\rightarrow$ 44,9; 768 $\rightarrow$ 45,4; 2048 $\rightarrow$ 45,6: tăng kênh ít tác dụng \\
Độ bền với ảnh hỏng (Bảng 5, Cityscapes-C) & SegFormer-B5 so với DeepLabv3+ (ResNet-50): ảnh gốc 82,4 / 76,6; nhiễu Gauss 57,8 / 6,5; làm mờ Gauss 69,8 / 57,7; nén JPEG 68,4 / 42,0 \\
\bottomrule
\end{tabularx}
\end{table}

\textbf{Kết luận của paper.} SegFormer đạt mIoU cao hơn các phương pháp cùng thời với ít tham số và FLOPs hơn; encoder không
positional encoding chịu được việc đổi độ phân giải lúc test; và mô hình bền hơn hẳn CNN khi ảnh bị nhiễu.

\subsection{Hạn chế của paper}
\begin{itemize}
  \item \st{So sánh độ bền không cùng cỡ.} Bảng 5 chỉ thử SegFormer-B5 (khoảng 85 triệu tham số) với các mô hình nhỏ hơn
        nhiều; paper không cho biết B2 có còn bền như vậy không, cũng không so với một CNN cùng cỡ được huấn luyện cùng
        công thức. Thí nghiệm ở Mục~\ref{sec:robust} của nhóm đo đúng điểm này.
  \item \st{Công thức huấn luyện khác nhau giữa các mô hình so sánh.} Kết quả của mô hình khác lấy lại từ paper gốc, với
        lịch huấn luyện và augmentation khác, nên chênh lệch không chỉ do kiến trúc.
  \item \st{Phụ thuộc dữ liệu và tính toán lớn.} Encoder cần pretrain ImageNet; mọi thí nghiệm dùng tập lớn (vài nghìn đến
        hàng trăm nghìn ảnh) và 8 GPU. Paper không thử chế độ ít nhãn.
  \item \st{Biên vật.} Bộ giải mã dự đoán ở 1/4 độ phân giải rồi phóng bilinear; với vật nhỏ hoặc biên mảnh (lông, tai) điều
        này có thể làm mất chi tiết. Paper chỉ minh hoạ định tính trên Cityscapes.
  \item Tác giả tự nêu: chưa rõ mô hình nhỏ nhất (3,7M tham số) có chạy được trên chip biên chỉ có 100~KB bộ nhớ không.
\end{itemize}

\subsection{Điều kiện trên tập của nhóm khác với paper}
Oxford-IIIT Pet khác các tập của paper ở ba điểm: (1) mỗi ảnh chỉ có \st{một con vật và nền}, nên ngữ cảnh toàn cảnh mà
attention mạnh ở ADE20K ít quan trọng hơn; (2) phần khó là \st{phân biệt 37 giống rất giống nhau}, gần với phân loại tinh
hơn là phân đoạn cảnh; (3) tập train chỉ có 3.303 ảnh, nhỏ hơn ADE20K khoảng 6 lần. Nhóm cũng huấn luyện ngắn hơn nhiều
(15 epoch, khoảng 3 nghìn vòng lặp) ở $384\times384$ trên một GPU T4.
